\section{Tiêu chuẩn đánh giá, hệ thống minh chứng và thực hành tự đánh giá}

\subsection{Cơ sở pháp lý và nguyên tắc phương pháp luận thiết lập tiêu chuẩn}

Căn cứ pháp lý trực tiếp để thiết lập hệ thống tiêu chuẩn đánh giá chất lượng tại Trường Cao đẳng Kỹ Nghệ II là Thông tư số 14/2024/TT-BLĐTBXH ngày 31 tháng 12 năm 2024 của Bộ trưởng Bộ Lao động -- Thương binh và Xã hội quy định tiêu chí, tiêu chuẩn kiểm định chất lượng giáo dục nghề nghiệp. Văn bản quy phạm pháp luật này đã xác lập một bước chuyển quan trọng trong tư duy quản trị chất lượng, chuyển từ việc kiểm tra hồ sơ giấy tờ thuần túy sang đánh giá năng lực thực chất, mức độ đạt chuẩn đầu ra và khả năng thích ứng của người học với thị trường lao động.

Về phương diện phương pháp luận, hệ thống tiêu chí, công cụ đo lường và khung minh chứng của nhà trường được xây dựng dựa trên bốn nguyên tắc quản trị chất lượng hiện đại. 

Trước hết là nguyên tắc bảo đảm tính tương thích pháp lý tuyệt đối. Toàn bộ cấu trúc đánh giá bám sát chặt chẽ tám tiêu chí và năm mươi tiêu chuẩn được quy định tại Phụ lục 01 ban hành kèm theo Thông tư số 14/2024/TT-BLĐTBXH, bảo đảm tính thống nhất và đồng bộ với hệ thống kiểm định chất lượng giáo dục nghề nghiệp quốc gia.

Thứ hai là nguyên tắc bảo đảm tính tích hợp tác nghiệp trong quản trị nội bộ. Mỗi tiêu chuẩn kiểm định của Thông tư số 14/2024/TT-BLĐTBXH không tồn tại như một yêu cầu hành chính biệt lập, mà được kết nối và ánh xạ trực tiếp với ít nhất một trong sáu mươi lăm quy trình vận hành nội bộ đã được Hiệu trưởng nhà trường ký ban hành. Cách tiếp cận này giúp hoạt động tự đánh giá trở thành một phần tự nhiên của công tác điều hành hàng ngày thay vì tạo thêm gánh nặng thủ tục khi đến kỳ kiểm định.

Thứ ba là nguyên tắc bảo đảm tính đo lường được của các chỉ báo chất lượng. Nhà trường chủ động cụ thể hóa các tiêu chuẩn mang tính định tính thành hệ thống các chỉ số đo lường hiệu năng cốt lõi có thể định lượng khách quan. Việc định lượng hóa các chỉ số giúp các cấp quản lý từ Ban Giám hiệu, Trưởng khoa cho đến giảng viên có thể theo dõi tiến độ thực hiện và đối soát dữ liệu một cách minh bạch, chính xác.

Thứ tư là nguyên tắc tuân thủ nghiêm ngặt bốn thuộc tính của minh chứng kiểm định. Mọi dữ liệu và hồ sơ được sử dụng làm minh chứng phục vụ tự đánh giá và đánh giá ngoài đều phải thỏa mãn đồng thời bốn yêu cầu cốt lõi: tính đầy đủ nhằm bao quát toàn diện các khía cạnh của tiêu chuẩn; tính tường minh thể hiện qua nguồn gốc pháp lý, thẩm quyền ban hành, ngày tháng, chữ ký và con dấu rõ ràng; tính tin cậy bảo đảm dữ liệu có thể kiểm chứng độc lập mà không bị sai lệch; và tính tương thích phản ánh đúng nội hàm, bản chất của chỉ báo đánh giá.

\subsection{Khung tiêu chuẩn và ma trận tích hợp minh chứng quản trị nội bộ}

Nhằm thiết lập cầu nối vững chắc giữa bộ tiêu chuẩn kiểm định quốc gia với thực tiễn quản trị tại cơ sở, Trường Cao đẳng Kỹ Nghệ II đã xây dựng Ma trận tích hợp giữa tám tiêu chí kiểm định theo Thông tư số 14/2024/TT-BLĐTBXH, danh mục minh chứng cốt lõi và mã quy trình nội bộ tương ứng của nhà trường.

\begin{small}
\setlength{\tabcolsep}{3.5pt}
\setstretch{1.15}
\begin{longtable}{C{1.1cm} L{3.0cm} C{0.8cm} L{3.6cm} L{2.5cm} L{2.5cm}}
\caption{\textbf{Ma trận tích hợp Bộ tiêu chuẩn kiểm định và 65 Quy trình nội bộ của Trường Cao đẳng Kỹ Nghệ II}} \label{tab:matran65qt} \\
\toprule
\textbf{Mã TC} & \textbf{Tên tiêu chí (Số tiêu chuẩn)} & \textbf{Điểm} & \textbf{Danh mục minh chứng cốt lõi} & \textbf{Quy trình nội bộ} & \textbf{Đơn vị chủ trì} \\
\midrule
\endfirsthead

\multicolumn{6}{c}{\textit{Bảng \ref{tab:matran65qt} (tiếp theo)}} \\[0.2cm]
\toprule
\textbf{Mã TC} & \textbf{Tên tiêu chí (Số tiêu chuẩn)} & \textbf{Điểm} & \textbf{Danh mục minh chứng cốt lõi} & \textbf{Quy trình nội bộ} & \textbf{Đơn vị chủ trì} \\
\midrule
\endhead

\midrule
\multicolumn{6}{r}{\textit{(Xem tiếp trang sau)}} \\
\endfoot

\bottomrule
\endlastfoot

TC-01 & Sứ mạng, mục tiêu, tổ chức và quản lý (5 tiêu chuẩn) & 10 & Chiến lược phát triển trường; Quy chế tổ chức và hoạt động; Quyết định thành lập Hội đồng trường; Sổ tay chất lượng; Báo cáo đánh giá nội bộ định kỳ. & QT/TCHCQT/01, QT/BĐCL/06, QT/BĐCL/07. & Phòng Tổ chức -- Hành chính -- Quản trị, Phòng Đảm bảo chất lượng \\
\addlinespace
TC-02 & Hoạt động đào tạo (8 tiêu chuẩn) & 16 & Đề án tuyển sinh hàng năm; Kế hoạch tiến độ đào tạo; Thời khóa biểu số hóa; Sổ theo dõi giảng dạy; Ngân hàng đề thi được phê duyệt. & ĐT/QT01/KHĐT, ĐT/QT02/TĐ, ĐT/QT03/TKB, ĐT/QT19/GDTT. & Phòng Đào tạo, các Khoa chuyên môn \\
\addlinespace
TC-03 & Nhà giáo, cán bộ quản lý, viên chức và người lao động (7 tiêu chuẩn) & 14 & Hồ sơ lý lịch khoa học; Văn bằng chuyên môn; Chứng chỉ sư phạm dạy nghề; Chứng chỉ kỹ năng nghề; Kế hoạch thực tế doanh nghiệp. & QT/TCHCQT/02, QT/TCHCQT/03, ĐT/QT10/DGG. & Phòng Tổ chức -- Hành chính -- Quản trị, các Khoa chuyên môn \\
\addlinespace
TC-04 & Chương trình đào tạo và giáo trình (7 tiêu chuẩn) & 14 & Chương trình chi tiết các ngành; Chuẩn đầu ra công bố công khai; Biên bản lấy ý kiến doanh nghiệp; Quyết định ban hành giáo trình. & ĐT/QT04/CT-GT, ĐT/QT09/GTĐT, QT/TCHCQT/11. & Phòng Đào tạo, Hội đồng Khoa học và Đào tạo, Thư viện \\
\addlinespace
TC-05 & Cơ sở vật chất, thiết bị đào tạo và thư viện (8 tiêu chuẩn) & 16 & Giấy chứng nhận quyền sử dụng đất; Danh mục thiết bị thực hành xưởng; Sổ bảo dưỡng máy móc định kỳ; Biên bản kiểm định an toàn; Tài liệu số hóa. & QT/TCHCQT/06, QT/TCHCQT/07, QT/TCHCQT/08. & Trung tâm Quản trị thiết bị, Thư viện, các Xưởng thực hành \\
\addlinespace
TC-06 & Nghiên cứu khoa học, chuyển giao và hợp tác quốc tế (4 tiêu chuẩn) & 8 & Thuyết minh và biên bản nghiệm thu đề tài nghiên cứu; Hợp đồng chuyển giao công nghệ; Danh mục bài báo; Biên bản ghi nhớ quốc tế. & QT/KHHTQT/01, QT/KHHTQT/02, QT/KHHTQT/04. & Ban Quản lý Khoa học và Hợp tác Quốc tế \\
\addlinespace
TC-07 & Dịch vụ người học (5 tiêu chuẩn) & 10 & Quyết định miễn giảm học phí; Danh sách chi trả học bổng qua ngân hàng; Sổ theo dõi cố vấn học tập; Biên bản đối thoại người học. & ĐT/QT28/HB, ĐT/QT31/HP, ĐT/QT34/CVHT. & Phòng Công tác sinh viên, Đoàn Thanh niên, Giáo viên chủ nhiệm \\
\addlinespace
TC-08 & Giám sát, đánh giá chất lượng (6 tiêu chuẩn) & 12 & Báo cáo tỷ lệ thôi học, tốt nghiệp; Báo cáo điều tra việc làm; Báo cáo khảo sát sự hài lòng; Kế hoạch hành động cải tiến sau tự đánh giá. & QT/BĐCL/01, QT/BĐCL/02, QT/BĐCL/04, ĐT/QT25/KSVL. & Phòng Đảm bảo chất lượng, Phòng Đào tạo, Trung tâm Quan hệ doanh nghiệp \\
\midrule
\textbf{Tổng} & \textbf{50 Tiêu chuẩn} & \textbf{100} & \textbf{Hệ thống minh chứng chuẩn hóa} & \textbf{Tích hợp 65 quy trình} & \textbf{Toàn trường} \\
\end{longtable}
\end{small}

\subsection{Bộ công cụ đo lường và đánh giá chất lượng đào tạo}

Để khắc phục khoảng trống thiếu công cụ đo lường thực tế và độ trễ thông tin trong thu thập phản hồi, Trường Cao đẳng Kỹ Nghệ II đã thiết kế hai công cụ đánh giá nòng cốt, bao gồm Phiếu khảo sát phản hồi người học theo thang đo Likert năm mức độ và Rubric đánh giá kỹ năng nghề nghiệp tích hợp năng lực số.

\subsubsection{Công cụ khảo sát ý kiến phản hồi của người học}

Công cụ khảo sát phản hồi người học có mã hiệu CC-BĐCL-01 được thiết kế dành riêng cho đối tượng học sinh, sinh viên các lớp cao đẳng chính quy. Thời điểm triển khai khảo sát được ấn định vào tuần thứ mười bốn của mỗi học kỳ, ngay trước khi bước vào kỳ thi kết thúc học phần. Để bảo đảm tính khách quan, trung thực và loại bỏ tâm lý e ngại của người học, quy trình khảo sát được thực hiện trực tuyến một trăm phần trăm qua ứng dụng di động nội bộ của nhà trường, tự động mã hóa danh tính sinh viên. Hệ thống sử dụng thang đo Likert năm mức độ được định nghĩa rõ ràng: Mức 1 là hoàn toàn không đồng ý, Mức 2 là không đồng ý, Mức 3 là phân vân hoặc bình thường, Mức 4 là đồng ý, và Mức 5 là hoàn toàn đồng ý.

Nội dung bảng hỏi khảo sát được chuẩn hóa chi tiết như sau:

\begin{quote}
\centering
\textbf{BẢNG HỎI KHẢO SÁT CHẤT LƯỢNG GIẢNG DẠY VÀ DỊCH VỤ ĐÀO TẠO}\\[0.1cm]
\textit{(Ban hành kèm theo Quy trình QT/BĐCL/02 -- Bản chuẩn hóa năm 2025)}\\[0.3cm]
\raggedright

\textbf{Phần A: Thông tin chung}\\[0.15cm]
Hệ thống tự động ghi nhận mã lớp học phần, tuyệt đối không lưu danh tính sinh viên tham gia khảo sát. Khóa đào tạo: \dots\dots; Khoa chuyên môn: \dots\dots; Môn học/mô-đun: \dots\dots; Giảng viên: \dots\dots\\[0.3cm]

\textbf{Phần B: Đánh giá nội dung và phương pháp giảng dạy của giảng viên}\\[0.15cm]
Người học lựa chọn mức độ phù hợp từ 1 đến 5 cho từng chỉ báo:
\begin{enumerate}
    \item Giảng viên phổ biến rõ ràng đề cương chi tiết, chuẩn đầu ra môn học và phương pháp đánh giá kết quả ngay trong buổi học đầu tiên.
    \item Giảng viên lên lớp đúng giờ, thực hiện đầy đủ thời lượng và nội dung giảng dạy theo đúng kế hoạch tiến độ đào tạo đã công bố.
    \item Phương pháp giảng dạy tích cực, kết hợp nhuần nhuyễn giữa lý thuyết cơ sở với thao tác mẫu thực hành trực quan và gợi mở tư duy người học.
    \item Giảng viên làm chủ công nghệ, sử dụng thành thạo phần mềm mô phỏng chuyên ngành và các thiết bị kỹ thuật số hỗ trợ bài giảng.
    \item Giảng viên tận tình hướng dẫn, uốn nắn chuẩn xác từng thao tác kỹ năng nghề và giải đáp kịp thời các thắc mắc của sinh viên tại xưởng máy.
    \item Việc đánh giá kết quả học tập bảo đảm công bằng, khách quan, công khai kịp thời điểm kiểm tra định kỳ và sửa bài chi tiết cho người học.
\end{enumerate}

\textbf{Phần C: Đánh giá cơ sở vật chất và dịch vụ hỗ trợ đào tạo}\\[0.15cm]
\begin{enumerate}
    \setcounter{enumi}{6}
    \item Phòng học lý thuyết và xưởng thực hành bảo đảm đủ ánh sáng, thông thoáng và đáp ứng đầy đủ các quy định về an toàn, vệ sinh lao động.
    \item Máy móc, thiết bị, dụng cụ thực hành tại xưởng hoạt động ổn định; vật tư thực hành tiêu hao được cấp phát đầy đủ, kịp thời theo bài học.
    \item Giáo trình, tài liệu tham khảo số hóa và hạ tầng mạng internet không dây tại xưởng đáp ứng tốt nhu cầu tra cứu và thực hành bài giảng.
    \item Đội ngũ cán bộ quản lý xưởng, giáo viên chủ nhiệm và chuyên viên hỗ trợ sinh viên có thái độ hòa nhã, tận tâm và giải quyết kịp thời các kiến nghị.
\end{enumerate}

\textbf{Phần D: Ý kiến đóng góp tự do và đề xuất cải tiến}\\[0.15cm]
Người học tự do chia sẻ cảm nhận về điều hài lòng nhất đối với môn học, đồng thời đưa ra các đề xuất cụ thể nhằm giúp nhà trường cải tiến chất lượng giảng dạy và dịch vụ xưởng trong các học kỳ tiếp theo.
\end{quote}

\subsubsection{Tiêu chí định lượng đánh giá kỹ năng nghề tích hợp năng lực số}

Công cụ Rubric đánh giá có mã hiệu RB-BĐCL-02 được thiết kế nhằm phục vụ đánh giá bài thi kết thúc mô-đun kỹ thuật tích hợp, điển hình như mô-đun Lắp đặt và lập trình điều khiển tự động hóa bằng bộ điều khiển lập trình (PLC) thuộc khối ngành Điện công nghiệp và Cơ điện tử. Thang điểm đánh giá gồm 10 điểm, phân thành bốn cấp độ năng lực tương ứng: Mức Kém (dưới 5,0 điểm), Mức Trung bình (từ 5,0 đến 6,9 điểm), Mức Khá (từ 7,0 đến 8,4 điểm) và Mức Xuất sắc (từ 8,5 đến 10,0 điểm).

\begin{table}[htbp]
\centering
\footnotesize
\setlength{\tabcolsep}{3.5pt}
\setstretch{1.15}
\caption{\textbf{Rubric đánh giá kỹ năng nghề nghiệp tích hợp năng lực số}}
\vspace{0.2cm}
\begin{tabularx}{\textwidth}{L{2.8cm} X X X X}
\toprule
\textbf{Tiêu chí và Trọng số} & \textbf{Mức Kém (< 50\%)} & \textbf{Trung bình (50--69\%)} & \textbf{Mức Khá (70--84\%)} & \textbf{Xuất sắc (85--100\%)} \\
\midrule
\textbf{1. An toàn và Chuẩn bị (15\%)} & Vi phạm quy tắc an toàn xưởng; thiếu bảo hộ lao động; bố trí bàn thực hành mất an toàn. & Có bảo hộ nhưng thao tác còn lúng túng; nhận biết rủi ro cơ bản; chuẩn bị chưa thật ngăn nắp. & Tuân thủ tốt quy định an toàn; trang bị bảo hộ đầy đủ; sắp xếp bàn xưởng ngăn nắp theo nguyên tắc 5S. & Tuân thủ tuyệt đối an toàn; chủ động rà soát an toàn cho bạn học xung quanh; chuẩn bị chuyên nghiệp. \\
\addlinespace
\textbf{2. Đấu nối phần cứng (25\%)} & Đấu nối sai sơ đồ nguyên lý; mối nối lỏng lẻo; gây chập mạch hoặc làm hỏng linh kiện thiết bị. & Đấu nối đúng sơ đồ cơ bản nhưng dây chưa ngay ngắn; cần giảng viên nhắc nhở lỗi tiếp xúc. & Đấu nối chính xác, chắc chắn, luồn dây gọn gàng; tự kiểm tra thông mạch độc lập trước khi đóng điện. & Đấu nối hoàn hảo, đạt chuẩn thẩm mỹ công nghiệp; đánh số đầu dây chuyên nghiệp; thao tác chuẩn xác. \\
\addlinespace
\textbf{3. Lập trình và Mô phỏng (35\%)} & Không xây dựng được thuật toán; không sử dụng được phần mềm; không mô phỏng được hệ thống. & Viết được đoạn mã lệnh cơ bản; thuật toán thô sơ; mô phỏng được một phần, cần hỗ trợ khi nạp mã. & Sử dụng thành thạo phần mềm; lập trình cấu trúc thuật toán mạch lạc, tối ưu; mô phỏng đạt yêu cầu. & Thuật toán thông minh; tích hợp giao diện giám sát số hóa người -- máy; tự kiểm tra và xử lý lỗi hoàn hảo. \\
\addlinespace
\textbf{4. Vận hành và Xử lý sự cố (15\%)} & Hệ thống không hoạt động; không xác định được nguyên nhân sự cố; thao tác hoảng loạn khi có lỗi. & Vận hành được ở chế độ bằng tay; khi phát sinh sự cố kỹ thuật mất nhiều thời gian tìm nguyên nhân. & Vận hành trơn tru cả ở chế độ tự động và bằng tay; khắc phục nhanh sự cố kỹ thuật thông thường. & Vận hành tối ưu; phân tích chính xác dữ liệu vận hành kỹ thuật; làm chủ hoàn toàn chu kỳ làm việc của máy. \\
\addlinespace
\textbf{5. Tác phong và Báo cáo số (10\%)} & Chậm tiến độ quy định; không ghi chép nhật ký kỹ thuật; thái độ làm việc thiếu tập trung. & Hoàn thành sát thời hạn; ghi nhật ký thực hành trên giấy còn sơ sài; vệ sinh xưởng chưa sạch sẽ. & Hoàn thành đúng tiến độ; xuất dữ liệu báo cáo kỹ thuật định dạng số rõ ràng; vệ sinh xưởng ngăn nắp. & Hoàn thành trước hạn; báo cáo số chuyên nghiệp có đồ thị trực quan; thể hiện đĩnh đạc tác phong công nghiệp. \\
\bottomrule
\end{tabularx}
\end{table}

\subsection{Báo cáo tự đánh giá thực nghiệm đối với tiêu chuẩn biến động người học}

Nhằm chứng minh năng lực thực hành tự đánh giá gắn liền với tư duy phản biện học thuật và đề xuất các giải pháp khắc phục triệt để khoảng trống chất lượng, dưới đây là Báo cáo tự đánh giá chi tiết cho Tiêu chuẩn 8.1 thuộc Tiêu chí 8 theo đúng cấu trúc chuẩn của Thông tư số 14/2024/TT-BLĐTBXH.

\subsubsection{Tiêu chuẩn 8.1: Theo dõi, đánh giá tỷ lệ người học bỏ học, thôi học, chuyển trường, lưu ban, tốt nghiệp}

\noindent\textit{(Nội dung tự đánh giá thuộc Tiêu chí 8: Giám sát, đánh giá chất lượng theo Thông tư số 14/2024/TT-BLĐTBXH)}\par\vspace{0.2cm}

\paragraph{1. Mô tả hiện trạng thực tế}
Trong giai đoạn từ năm 2022 đến năm 2025, Trường Cao đẳng Kỹ Nghệ II đã ban hành và duy trì hệ thống văn bản quản lý liên quan đến công tác quản lý tiến độ đào tạo và biến động người học, trọng tâm là Quy chế đào tạo ban hành theo Quyết định số 188/QĐ-CĐKNII và Quy trình quản lý tiến độ đào tạo có mã hiệu ĐT/QT02/TĐ. Định kỳ từng học kỳ, Phòng Đào tạo phối hợp cùng Phòng Công tác sinh viên và các Khoa chuyên môn trích xuất dữ liệu từ phần mềm quản lý đào tạo để lập danh sách theo dõi học sinh, sinh viên nghỉ học, ban hành các thông báo cảnh báo học vụ và trình Hiệu trưởng ký các quyết định buộc thôi học hoặc công nhận tốt nghiệp theo đúng thẩm quyền được minh chứng tại tài liệu có mã hiệu H8.08.01.01.

Dữ liệu thống kê thực tế về biến động người học trong ba khóa đào tạo cao đẳng chính quy gần nhất của nhà trường được phản ánh chi tiết trong bảng dưới đây:

\begin{table}[htbp]
\centering
\small
\setlength{\tabcolsep}{3.5pt}
\setstretch{1.2}
\caption{\textbf{Thống kê tỷ lệ duy trì sĩ số và biến động người học tại Trường Cao đẳng Kỹ Nghệ II}}
\vspace{0.2cm}
\begin{tabularx}{\textwidth}{L{2.8cm} C{1.5cm} C{1.6cm} C{1.5cm} C{1.5cm} C{1.5cm} X}
\toprule
\textbf{Khóa đào tạo} & \textbf{Quy mô đầu vào} & \textbf{Tốt nghiệp đúng hạn} & \textbf{Tỷ lệ tốt nghiệp} & \textbf{Thôi học lũy kế} & \textbf{Tỷ lệ hao hụt} & \textbf{Chỉ tiêu mục tiêu chất lượng} \\
\midrule
\textbf{Khóa K45 (2021 -- 2024)} & 1.850 & 1.425 & 77,0\% & 425 & 23,0\% & Duy trì sĩ số $\ge$ 75\% \\
\addlinespace
\textbf{Khóa K46 (2022 -- 2025)} & 2.100 & 1.638 & 78,0\% & 462 & 22,0\% & Duy trì sĩ số $\ge$ 75\% \\
\addlinespace
\textbf{Khóa K47 (2023 -- 2026)} & 2.250 & Đang đào tạo & Dự kiến > 78\% & 382 & 17,0\% & Duy trì sĩ số $\ge$ 75\% \\
\bottomrule
\end{tabularx}
\end{table}

Dữ liệu thực tế cho thấy tỷ lệ duy trì sĩ số bình quân của Trường Cao đẳng Kỹ Nghệ II dao động từ 77,0\% đến 78,0\%, hoàn thành chỉ tiêu đề ra tại Mục tiêu chất lượng số 3 năm 2025 của trường là duy trì sĩ số học sinh, sinh viên từ 75\% trở lên được lưu trữ tại minh chứng H8.08.01.02. Công tác xét công nhận tốt nghiệp được thực hiện đúng quy chế đào tạo; toàn bộ sinh viên đủ điều kiện đều được cấp văn bằng tốt nghiệp đúng thời hạn quy định theo minh chứng H8.08.01.03.

\paragraph{2. Điểm mạnh}
Nhà trường đã thiết lập được quy trình phối hợp tương đối nền nếp trong việc thống kê và đối soát dữ liệu biến động sĩ số giữa Phòng Đào tạo, Phòng Công tác sinh viên và các đơn vị chuyên môn. Tỷ lệ sinh viên tốt nghiệp đúng hạn so với tổng quy mô nhập học ban đầu đạt xấp xỉ 78\%, phản ánh sự cố gắng trong khâu tổ chức đào tạo và giảng dạy. Đồng thời, toàn bộ hồ sơ công nhận tốt nghiệp và cấp phát văn bằng chứng chỉ được lưu trữ chặt chẽ và quản lý trên phần mềm số hóa, bảo đảm tính pháp lý và độ chính xác cao.

\paragraph{3. Điểm tồn tại, hạn chế}
Bên cạnh các kết quả đạt được, công tác quản trị biến động người học tại trường bộc lộ ba tồn tại căn bản. Tồn tại thứ nhất là tư duy quản trị chấp nhận mức hao hụt sĩ số quá cao. Việc nhà trường đặt chỉ tiêu duy trì sĩ số từ 75\% trở lên đã vô hình trung tạo ra tâm lý thỏa mãn của các đơn vị khi tỷ lệ bỏ học thực tế từ 22\% đến 23\% vẫn được coi là hoàn thành nhiệm vụ. Đối với một cơ sở giáo dục nghề nghiệp tự chủ tài chính toàn diện, việc để mất đi gần một phần tư quy mô người học trong mỗi khóa đào tạo gây ra lãng phí nghiêm trọng nguồn lực giảng dạy và làm sụt giảm nguồn thu bền vững. Tồn tại thứ hai là phương thức quản lý biến động người học mang tính bị động và xử lý sự việc đã rồi. Nhà trường hiện chỉ thực hiện thống kê cơ học và ban hành quyết định xóa tên khi sinh viên đã bỏ học tự do quá số buổi quy định hoặc nợ học phí kéo dài mà hoàn toàn thiếu cơ chế tìm hiểu nguyên nhân gốc rễ để can thiệp kịp thời. Tồn tại thứ ba là sự vắng bóng của một công cụ cảnh báo sớm tự động dựa trên phân tích dữ liệu chuyên cần và điểm kiểm tra quá trình, dẫn đến việc nhà trường bỏ lỡ thời điểm vàng để hỗ trợ người học trong năm thứ nhất, giai đoạn chiếm tới 70\% tổng số sinh viên thôi học toàn khóa.

\paragraph{4. Kế hoạch hành động cải tiến chất lượng}
Để giải quyết triệt để các tồn tại nêu trên và nâng cao tỷ lệ duy trì người học, Trường Cao đẳng Kỹ Nghệ II xác lập Kế hoạch hành động cải tiến chất lượng chi tiết cho giai đoạn 2026 -- 2027 được trình bày trong bảng dưới đây:

\begin{table}[htbp]
\centering
\small
\setlength{\tabcolsep}{4pt}
\setstretch{1.2}
\caption{\textbf{Kế hoạch hành động cải tiến Tiêu chuẩn 8.1}}
\vspace{0.2cm}
\begin{tabularx}{\textwidth}{L{3.8cm} X L{3.0cm} L{2.2cm} C{1.8cm}}
\toprule
\textbf{Nội dung công việc} & \textbf{Mục tiêu cần đạt} & \textbf{Đơn vị chủ trì và phối hợp} & \textbf{Nguồn lực} & \textbf{Thời hạn} \\
\midrule
\textbf{1. Nâng chuẩn chỉ tiêu duy trì sĩ số} trong Bảng Mục tiêu chất lượng 2026. & Nâng tỷ lệ duy trì sĩ số tối thiểu từ 75\% lên 88\%, khống chế bỏ học dưới 12\%. & Hiệu trưởng, Phòng Đảm bảo chất lượng & Văn bản pháp lý nội bộ & Tháng 01/2026 \\
\addlinespace
\textbf{2. Xây dựng Hệ thống cảnh báo sớm} trên phần mềm quản trị đào tạo. & Tự động gửi cảnh báo khi vắng 3 buổi liên tiếp hoặc điểm quá trình dưới 4,0. & Phòng Đào tạo, Trung tâm Công nghệ thông tin, Phòng Đảm bảo chất lượng & Kinh phí chuyển đổi số (120 triệu) & Tháng 06/2026 \\
\addlinespace
\textbf{3. Thiết lập mạng lưới can thiệp} hỗ trợ sinh viên có nguy cơ thôi học. & 100\% sinh viên nhận cảnh báo học vụ đều được giáo viên và Đoàn trường hỗ trợ. & Phòng Công tác sinh viên, Đoàn Thanh niên, các Khoa & Quỹ Hỗ trợ người học và học bổng & Từ học kỳ 1 năm 2026 \\
\addlinespace
\textbf{4. Ban hành Quy trình thao tác chuẩn} theo dõi và can thiệp chủ động. & Chuẩn hóa quy trình xử lý thông tin biến động sĩ số thành quy trình chính thức. & Phòng Đảm bảo chất lượng, Phòng Đào tạo & Hệ thống tài liệu nội bộ & Tháng 09/2026 \\
\bottomrule
\end{tabularx}
\end{table}

\paragraph{5. Tự đánh giá tiêu chuẩn}
Căn cứ theo các tiêu chí và chỉ báo quy định tại Thông tư số 14/2024/TT-BLĐTBXH, nhà trường tự đánh giá Tiêu chuẩn 8.1 đạt yêu cầu với mức điểm 2,0 trên thang điểm 2,0. Mặc dù đạt điểm tối đa trên hồ sơ kiểm định, nhà trường chủ động xếp tiêu chuẩn này vào danh mục ưu tiên cải tiến cấp độ một trong kế hoạch quản trị rủi ro giai đoạn 2026 -- 2030 nhằm hướng tới sự phát triển chất lượng thực chất và bền vững. Điểm tự đánh giá: 2,0 / 2,0 điểm.

\paragraph{6. Danh mục minh chứng kèm theo Tiêu chuẩn 8.1}
Danh mục minh chứng xác thực đính kèm báo cáo tự đánh giá tiêu chuẩn bao gồm: tài liệu mã hiệu H8.08.01.01 là Quyết định ban hành Quy chế đào tạo và Quy trình quản lý tiến độ đào tạo có mã hiệu ĐT/QT02/TĐ; tài liệu mã hiệu H8.08.01.02 là Quyết định ban hành Mục tiêu chất lượng năm 2024 và năm 2025 của Hiệu trưởng; tài liệu mã hiệu H8.08.01.03 là Báo cáo thống kê số liệu sinh viên tốt nghiệp và các Quyết định công nhận tốt nghiệp giai đoạn 2022 -- 2025; và tài liệu mã hiệu H8.08.01.04 là Tập hợp các Quyết định cảnh báo học vụ, Quyết định xóa tên học sinh, sinh viên thôi học, bỏ học giai đoạn 2022 -- 2025.
